% --- Archon physics formalization source begin ---
% archon:physics
% archon:covers PhyXMiniProblems/problem_phyx_mini_0534.lean
% archon:source-report reports/phyx_mini/problem_phyx_mini_0534.source.json

\chapter{Physics problem phyx\_mini\_0534}
\label{ch:PhyXMiniProblems_problem_phyx_mini_0534}

\paragraph{Problem source.}
\#\# Physical scenario

Particles incident from the left in figure are confronted with a step in potential energy. The step has a height \textbackslash{}( U \textbackslash{}) at \textbackslash{}( x = 0 \textbackslash{}). The particles have energy \textbackslash{}( E > U \textbackslash{}). Classically, all the particles would continue moving forward with reduced speed. According to quantum mechanics, however, a fraction of the particles are reflected at the step.

\#\# Question

Find its probability of being transmitted.

\#\# Answer choices

- A: A: 0.939

- B: B: 0.866

- C: C: 0.745

- D: D: 0.908

\#\# Figure caption (auxiliary; use the image as primary evidence)

The image depicts a diagram related to particle physics or quantum mechanics. It shows an "Incoming particle" represented by a red dot with an arrow pointing to the right. The path of the particle is horizontal. There are two levels in the diagram, marked by horizontal lines. 

- The lower horizontal line is labeled "U = 0," indicating a region of zero potential energy.
- There is a vertical step up to the second level, which indicates a potential energy barrier.
- The height of this step is labeled "U," representing the potential energy, and a vertical arrow points from the lower line to the top of the step.
- Another vertical arrow, labeled "E," points upward from the lower line, representing the total energy of the incoming particle.

The diagram illustrates a scenario where a particle with energy \textbackslash{}( E \textbackslash{}) approaches a potential energy step where the potential changes from 0 to \textbackslash{}( U \textbackslash{}).

\#\# Dataset metadata

Quantum Phenomena; Physical Model Grounding Reasoning, Numerical Reasoning

\paragraph{Recorded answer/context.}
D: D: 0.908

\paragraph{Figure/image path.}
/root/proposal\_for\_physic/science-mango/phyx\_mini\_run/phyx\_data/test\_image/534.png

\paragraph{Formalization target.}
create a compiling Lean file with sorry bodies at `PhyXMiniProblems/problem\_phyx\_mini\_0534.lean`. The Lean declarations must preserve the physical quantities, dimensions or dimensional roles, figure labels, governing-law hypotheses, and final relation expressed by this problem.
Use Mathlib/Physlib names found through LeanExplore where available. If a physics API is missing, introduce faithful local abstractions rather than scalar placeholder aliases.

\begin{theorem}[Physics formalization target]
\label{thm:physics:phyx_mini_0534:target}
\lean{PhyXMiniProblems.ProblemPhyXMini0534.transmissionProbabilityForFigurePotentialStep}
\uses{def:physics:phyx-mini-0534:phyxminiproblems-problemphyxmini0534-quantumpotentialstepsetup, def:physics:phyx-mini-0534:phyxminiproblems-problemphyxmini0534-transmissionprobability, def:physics:phyx-mini-0534:phyxminiproblems-problemphyxmini0534-matchespotentialstepscenario, def:physics:phyx-mini-0534:phyxminiproblems-problemphyxmini0534-matchessuppliedpotentialstepfigure, def:physics:phyx-mini-0534:phyxminiproblems-problemphyxmini0534-useslinearfigureenergyscale, def:physics:phyx-mini-0534:phyxminiproblems-problemphyxmini0534-hasphysicalpotentialstepparameters, def:physics:phyx-mini-0534:phyxminiproblems-problemphyxmini0534-usesstandardreducedplanckconstant, def:physics:phyx-mini-0534:phyxminiproblems-problemphyxmini0534-satisfiesquantumpotentialsteplaws, def:physics:phyx-mini-0534:phyxminiproblems-problemphyxmini0534-isnearestanswerchoice, def:physics:phyx-mini-0534:phyxminiproblems-problemphyxmini0534-agreeswhenroundedtothreedecimals}
The assigned autoformalize agent should translate this physics problem into Lean declarations in the covered file, with theorem and lemma proof bodies written as `by sorry`.
\end{theorem}
\begin{proof}
This is an autoformalization task, not a proof task. Produce statements that can later be proved by the physics prover without weakening the physical model.
\end{proof}

% --- Archon declaration topology begin ---
\section{Declaration topology}
\begin{definition}[Mass Quantity]
  \label{def:physics:phyx-mini-0534:phyxminiproblems-problemphyxmini0534-massquantity}
  \lean{PhyXMiniProblems.ProblemPhyXMini0534.MassQuantity}
  A nonnegative, unit-independent physical mass
\end{definition}
\begin{proof}
  This is the declared model definition of Mass Quantity.
\end{proof}
\begin{definition}[action Dimension]
  \label{def:physics:phyx-mini-0534:phyxminiproblems-problemphyxmini0534-actiondimension}
  \lean{PhyXMiniProblems.ProblemPhyXMini0534.actionDimension}
  The physical dimension of action, equivalently energy times time
\end{definition}
\begin{proof}
  This is the declared model definition of action Dimension.
\end{proof}
\begin{definition}[Action Quantity]
  \label{def:physics:phyx-mini-0534:phyxminiproblems-problemphyxmini0534-actionquantity}
  \lean{PhyXMiniProblems.ProblemPhyXMini0534.ActionQuantity}
  \uses{def:physics:phyx-mini-0534:phyxminiproblems-problemphyxmini0534-actiondimension}
  A nonnegative, unit-independent physical action
\end{definition}
\begin{proof}
  This is the declared model definition of Action Quantity.
\end{proof}
\begin{definition}[Signed Length Quantity]
  \label{def:physics:phyx-mini-0534:phyxminiproblems-problemphyxmini0534-signedlengthquantity}
  \lean{PhyXMiniProblems.ProblemPhyXMini0534.SignedLengthQuantity}
  A signed, unit-independent position along the scattering axis
\end{definition}
\begin{proof}
  This is the declared model definition of Signed Length Quantity.
\end{proof}
\begin{definition}[Wave Number Quantity]
  \label{def:physics:phyx-mini-0534:phyxminiproblems-problemphyxmini0534-wavenumberquantity}
  \lean{PhyXMiniProblems.ProblemPhyXMini0534.WaveNumberQuantity}
  A nonnegative wave-number magnitude, carrying inverse-length dimension
\end{definition}
\begin{proof}
  This is the declared model definition of Wave Number Quantity.
\end{proof}
\begin{definition}[energy In Joules]
  \label{def:physics:phyx-mini-0534:phyxminiproblems-problemphyxmini0534-energyinjoules}
  \lean{PhyXMiniProblems.ProblemPhyXMini0534.energyInJoules}
  Coherent-SI readout of an energy, in joules
\end{definition}
\begin{proof}
  This is the declared model definition of energy In Joules.
\end{proof}
\begin{definition}[mass In Kilograms]
  \label{def:physics:phyx-mini-0534:phyxminiproblems-problemphyxmini0534-massinkilograms}
  \lean{PhyXMiniProblems.ProblemPhyXMini0534.massInKilograms}
  \uses{def:physics:phyx-mini-0534:phyxminiproblems-problemphyxmini0534-massquantity}
  Coherent-SI readout of a mass, in kilograms
\end{definition}
\begin{proof}
  This is the declared model definition of mass In Kilograms.
\end{proof}
\begin{definition}[action In Joule Seconds]
  \label{def:physics:phyx-mini-0534:phyxminiproblems-problemphyxmini0534-actioninjouleseconds}
  \lean{PhyXMiniProblems.ProblemPhyXMini0534.actionInJouleSeconds}
  \uses{def:physics:phyx-mini-0534:phyxminiproblems-problemphyxmini0534-actionquantity}
  Coherent-SI readout of an action, in joule-seconds
\end{definition}
\begin{proof}
  This is the declared model definition of action In Joule Seconds.
\end{proof}
\begin{definition}[position In Meters]
  \label{def:physics:phyx-mini-0534:phyxminiproblems-problemphyxmini0534-positioninmeters}
  \lean{PhyXMiniProblems.ProblemPhyXMini0534.positionInMeters}
  \uses{def:physics:phyx-mini-0534:phyxminiproblems-problemphyxmini0534-signedlengthquantity}
  Read a signed position in metres
\end{definition}
\begin{proof}
  This is the declared model definition of position In Meters.
\end{proof}
\begin{definition}[wave Number In Inverse Meters]
  \label{def:physics:phyx-mini-0534:phyxminiproblems-problemphyxmini0534-wavenumberininversemeters}
  \lean{PhyXMiniProblems.ProblemPhyXMini0534.waveNumberInInverseMeters}
  \uses{def:physics:phyx-mini-0534:phyxminiproblems-problemphyxmini0534-wavenumberquantity}
  Read a wave-number magnitude in inverse metres
\end{definition}
\begin{proof}
  This is the declared model definition of wave Number In Inverse Meters.
\end{proof}
\begin{definition}[Potential Region]
  \label{def:physics:phyx-mini-0534:phyxminiproblems-problemphyxmini0534-potentialregion}
  \lean{PhyXMiniProblems.ProblemPhyXMini0534.PotentialRegion}
  The two constant-potential regions separated by x = 0
\end{definition}
\begin{proof}
  This is the declared model definition of Potential Region.
\end{proof}
\begin{definition}[Scattering Component]
  \label{def:physics:phyx-mini-0534:phyxminiproblems-problemphyxmini0534-scatteringcomponent}
  \lean{PhyXMiniProblems.ProblemPhyXMini0534.ScatteringComponent}
  The three plane-wave components in above-barrier step scattering
\end{definition}
\begin{proof}
  This is the declared model definition of Scattering Component.
\end{proof}
\begin{definition}[Horizontal Direction]
  \label{def:physics:phyx-mini-0534:phyxminiproblems-problemphyxmini0534-horizontaldirection}
  \lean{PhyXMiniProblems.ProblemPhyXMini0534.HorizontalDirection}
  The two directions along the horizontal scattering axis
\end{definition}
\begin{proof}
  This is the declared model definition of Horizontal Direction.
\end{proof}
\begin{definition}[Potential Step Figure]
  \label{def:physics:phyx-mini-0534:phyxminiproblems-problemphyxmini0534-potentialstepfigure}
  \lean{PhyXMiniProblems.ProblemPhyXMini0534.PotentialStepFigure}
  Visible contents and rounded measurements of the 428 × 236 source raster. The energy and step arrows have rounded tip-to-tip vertical spans of 154 and 110 pixels, respectively, giving the intended drawn scale E : U = 7 : 5. The pixel fields are dimensionless image readouts, not physical energies
\end{definition}
\begin{proof}
  This is the declared model definition of Potential Step Figure.
\end{proof}
\begin{definition}[Quantum Potential Step Setup]
  \label{def:physics:phyx-mini-0534:phyxminiproblems-problemphyxmini0534-quantumpotentialstepsetup}
  \lean{PhyXMiniProblems.ProblemPhyXMini0534.QuantumPotentialStepSetup}
  \uses{def:physics:phyx-mini-0534:phyxminiproblems-problemphyxmini0534-massquantity, def:physics:phyx-mini-0534:phyxminiproblems-problemphyxmini0534-actionquantity, def:physics:phyx-mini-0534:phyxminiproblems-problemphyxmini0534-signedlengthquantity, def:physics:phyx-mini-0534:phyxminiproblems-problemphyxmini0534-wavenumberquantity, def:physics:phyx-mini-0534:phyxminiproblems-problemphyxmini0534-potentialregion, def:physics:phyx-mini-0534:phyxminiproblems-problemphyxmini0534-scatteringcomponent, def:physics:phyx-mini-0534:phyxminiproblems-problemphyxmini0534-horizontaldirection, def:physics:phyx-mini-0534:phyxminiproblems-problemphyxmini0534-potentialstepfigure}
  Independent quantities in the scattering experiment. The two outcome probabilities and the regional wave numbers remain unknown fields here
\end{definition}
\begin{proof}
  This is the declared model definition of Quantum Potential Step Setup.
\end{proof}
\begin{definition}[transmission Probability]
  \label{def:physics:phyx-mini-0534:phyxminiproblems-problemphyxmini0534-transmissionprobability}
  \lean{PhyXMiniProblems.ProblemPhyXMini0534.transmissionProbability}
  \uses{def:physics:phyx-mini-0534:phyxminiproblems-problemphyxmini0534-quantumpotentialstepsetup}
  The probability requested by the problem
\end{definition}
\begin{proof}
  This is the declared model definition of transmission Probability.
\end{proof}
\begin{definition}[reflection Probability]
  \label{def:physics:phyx-mini-0534:phyxminiproblems-problemphyxmini0534-reflectionprobability}
  \lean{PhyXMiniProblems.ProblemPhyXMini0534.reflectionProbability}
  \uses{def:physics:phyx-mini-0534:phyxminiproblems-problemphyxmini0534-quantumpotentialstepsetup}
  The reflected fraction mentioned in the scenario
\end{definition}
\begin{proof}
  This is the declared model definition of reflection Probability.
\end{proof}
\begin{definition}[Matches Potential Step Scenario]
  \label{def:physics:phyx-mini-0534:phyxminiproblems-problemphyxmini0534-matchespotentialstepscenario}
  \lean{PhyXMiniProblems.ProblemPhyXMini0534.MatchesPotentialStepScenario}
  \uses{def:physics:phyx-mini-0534:phyxminiproblems-problemphyxmini0534-energyinjoules, def:physics:phyx-mini-0534:phyxminiproblems-problemphyxmini0534-positioninmeters, def:physics:phyx-mini-0534:phyxminiproblems-problemphyxmini0534-quantumpotentialstepsetup}
  The step is at the origin, the left potential is zero, and the right potential is the stated step height. The component roles encode incidence from the left, reflection back to the left, and transmission into the right region
\end{definition}
\begin{proof}
  This is the declared model definition of Matches Potential Step Scenario.
\end{proof}
\begin{definition}[Matches Supplied Potential Step Figure]
  \label{def:physics:phyx-mini-0534:phyxminiproblems-problemphyxmini0534-matchessuppliedpotentialstepfigure}
  \lean{PhyXMiniProblems.ProblemPhyXMini0534.MatchesSuppliedPotentialStepFigure}
  \uses{def:physics:phyx-mini-0534:phyxminiproblems-problemphyxmini0534-potentialstepfigure}
  Qualitative labels and rounded raster measurements from image 534.png. This premise contains neither a probability nor an answer-choice label
\end{definition}
\begin{proof}
  This is the declared model definition of Matches Supplied Potential Step Figure.
\end{proof}
\begin{definition}[Uses Linear Figure Energy Scale]
  \label{def:physics:phyx-mini-0534:phyxminiproblems-problemphyxmini0534-useslinearfigureenergyscale}
  \lean{PhyXMiniProblems.ProblemPhyXMini0534.UsesLinearFigureEnergyScale}
  \uses{def:physics:phyx-mini-0534:phyxminiproblems-problemphyxmini0534-energyinjoules, def:physics:phyx-mini-0534:phyxminiproblems-problemphyxmini0534-quantumpotentialstepsetup}
  The diagram uses one common linear vertical energy scale for its E and U arrows. Cross-multiplication avoids dividing by an image length. This is a figure-calibration law, not a transmission-probability formula
\end{definition}
\begin{proof}
  This is the declared model definition of Uses Linear Figure Energy Scale.
\end{proof}
\begin{definition}[Has Physical Potential Step Parameters]
  \label{def:physics:phyx-mini-0534:phyxminiproblems-problemphyxmini0534-hasphysicalpotentialstepparameters}
  \lean{PhyXMiniProblems.ProblemPhyXMini0534.HasPhysicalPotentialStepParameters}
  \uses{def:physics:phyx-mini-0534:phyxminiproblems-problemphyxmini0534-energyinjoules, def:physics:phyx-mini-0534:phyxminiproblems-problemphyxmini0534-massinkilograms, def:physics:phyx-mini-0534:phyxminiproblems-problemphyxmini0534-actioninjouleseconds, def:physics:phyx-mini-0534:phyxminiproblems-problemphyxmini0534-wavenumberininversemeters, def:physics:phyx-mini-0534:phyxminiproblems-problemphyxmini0534-quantumpotentialstepsetup, def:physics:phyx-mini-0534:phyxminiproblems-problemphyxmini0534-transmissionprobability, def:physics:phyx-mini-0534:phyxminiproblems-problemphyxmini0534-reflectionprobability}
  Positivity, above-step energy, and probability bounds for the physical branch
\end{definition}
\begin{proof}
  This is the declared model definition of Has Physical Potential Step Parameters.
\end{proof}
\begin{definition}[Uses Standard Reduced Planck Constant]
  \label{def:physics:phyx-mini-0534:phyxminiproblems-problemphyxmini0534-usesstandardreducedplanckconstant}
  \lean{PhyXMiniProblems.ProblemPhyXMini0534.UsesStandardReducedPlanckConstant}
  \uses{def:physics:phyx-mini-0534:phyxminiproblems-problemphyxmini0534-actioninjouleseconds, def:physics:phyx-mini-0534:phyxminiproblems-problemphyxmini0534-quantumpotentialstepsetup}
  Physlib's standard reduced Planck constant, calibrated in SI joule-seconds. This datum does not constrain either outcome probability
\end{definition}
\begin{proof}
  This is the declared model definition of Uses Standard Reduced Planck Constant.
\end{proof}
\begin{definition}[Satisfies Quantum Potential Step Laws]
  \label{def:physics:phyx-mini-0534:phyxminiproblems-problemphyxmini0534-satisfiesquantumpotentialsteplaws}
  \lean{PhyXMiniProblems.ProblemPhyXMini0534.SatisfiesQuantumPotentialStepLaws}
  \uses{def:physics:phyx-mini-0534:phyxminiproblems-problemphyxmini0534-energyinjoules, def:physics:phyx-mini-0534:phyxminiproblems-problemphyxmini0534-massinkilograms, def:physics:phyx-mini-0534:phyxminiproblems-problemphyxmini0534-actioninjouleseconds, def:physics:phyx-mini-0534:phyxminiproblems-problemphyxmini0534-wavenumberininversemeters, def:physics:phyx-mini-0534:phyxminiproblems-problemphyxmini0534-potentialregion, def:physics:phyx-mini-0534:phyxminiproblems-problemphyxmini0534-quantumpotentialstepsetup, def:physics:phyx-mini-0534:phyxminiproblems-problemphyxmini0534-transmissionprobability, def:physics:phyx-mini-0534:phyxminiproblems-problemphyxmini0534-reflectionprobability}
  The governing laws for a stationary one-dimensional upward step with E > U: the free-region Schrödinger dispersion relation ℏ² k² = 2m(E - V) on both sides; the probability-flux transmission and reflection coefficients obtained by matching the wave function and its derivative at the sharp interface; conservation of incident probability flux. All formulas are stated for the setup's still-unknown energies and wave numbers. None specializes the transmission probability to the current figure's numerical answer
\end{definition}
\begin{proof}
  This is the declared model definition of Satisfies Quantum Potential Step Laws.
\end{proof}
\begin{lemma}[transmitted To Incident Wave Number Ratio]
  \label{lem:physics:phyx-mini-0534:phyxminiproblems-problemphyxmini0534-transmittedtoincidentwavenumberratio}
  \lean{PhyXMiniProblems.ProblemPhyXMini0534.transmittedToIncidentWaveNumberRatio}
  \uses{def:physics:phyx-mini-0534:phyxminiproblems-problemphyxmini0534-wavenumberininversemeters, def:physics:phyx-mini-0534:phyxminiproblems-problemphyxmini0534-quantumpotentialstepsetup, def:physics:phyx-mini-0534:phyxminiproblems-problemphyxmini0534-matchespotentialstepscenario, def:physics:phyx-mini-0534:phyxminiproblems-problemphyxmini0534-matchessuppliedpotentialstepfigure, def:physics:phyx-mini-0534:phyxminiproblems-problemphyxmini0534-useslinearfigureenergyscale, def:physics:phyx-mini-0534:phyxminiproblems-problemphyxmini0534-hasphysicalpotentialstepparameters, def:physics:phyx-mini-0534:phyxminiproblems-problemphyxmini0534-satisfiesquantumpotentialsteplaws}
  The figure scale and regional dispersion imply k\_right / k\_left = sqrt (2/7). This is a derived intermediate result, not an assumption and not the requested transmission probability
\end{lemma}
\begin{proof}
  The conclusion follows from the governing relations and figure readouts named in the statement of transmitted To Incident Wave Number Ratio.
\end{proof}
\begin{definition}[Answer Choice]
  \label{def:physics:phyx-mini-0534:phyxminiproblems-problemphyxmini0534-answerchoice}
  \lean{PhyXMiniProblems.ProblemPhyXMini0534.AnswerChoice}
  Labels of the four transmission-probability choices
\end{definition}
\begin{proof}
  This is the declared model definition of Answer Choice.
\end{proof}
\begin{definition}[displayed Transmission Probability]
  \label{def:physics:phyx-mini-0534:phyxminiproblems-problemphyxmini0534-displayedtransmissionprobability}
  \lean{PhyXMiniProblems.ProblemPhyXMini0534.displayedTransmissionProbability}
  \uses{def:physics:phyx-mini-0534:phyxminiproblems-problemphyxmini0534-answerchoice}
  Dimensionless probability printed beside each answer choice
\end{definition}
\begin{proof}
  This is the declared model definition of displayed Transmission Probability.
\end{proof}
\begin{definition}[recorded Dataset Answer]
  \label{def:physics:phyx-mini-0534:phyxminiproblems-problemphyxmini0534-recordeddatasetanswer}
  \lean{PhyXMiniProblems.ProblemPhyXMini0534.recordedDatasetAnswer}
  \uses{def:physics:phyx-mini-0534:phyxminiproblems-problemphyxmini0534-answerchoice}
  Dataset metadata recording answer D; this is not a theorem premise
\end{definition}
\begin{proof}
  This is the declared model definition of recorded Dataset Answer.
\end{proof}
\begin{definition}[Is Nearest Answer Choice]
  \label{def:physics:phyx-mini-0534:phyxminiproblems-problemphyxmini0534-isnearestanswerchoice}
  \lean{PhyXMiniProblems.ProblemPhyXMini0534.IsNearestAnswerChoice}
  \uses{def:physics:phyx-mini-0534:phyxminiproblems-problemphyxmini0534-answerchoice, def:physics:phyx-mini-0534:phyxminiproblems-problemphyxmini0534-displayedtransmissionprobability}
  A displayed choice is at least as close as every alternative
\end{definition}
\begin{proof}
  This is the declared model definition of Is Nearest Answer Choice.
\end{proof}
\begin{definition}[Agrees When Rounded To Three Decimals]
  \label{def:physics:phyx-mini-0534:phyxminiproblems-problemphyxmini0534-agreeswhenroundedtothreedecimals}
  \lean{PhyXMiniProblems.ProblemPhyXMini0534.AgreesWhenRoundedToThreeDecimals}
  \uses{def:physics:phyx-mini-0534:phyxminiproblems-problemphyxmini0534-answerchoice, def:physics:phyx-mini-0534:phyxminiproblems-problemphyxmini0534-displayedtransmissionprobability}
  Agreement with a probability printed to three decimal places. Half of one unit in the final decimal place is 1 / 2000
\end{definition}
\begin{proof}
  This is the declared model definition of Agrees When Rounded To Three Decimals.
\end{proof}
% --- Archon declaration topology end ---
% --- Archon physics formalization source end ---
